\documentclass[11pt]{article}
\newcommand{\ListaTitulo}{Lista 06: Probabilidade condicional e independência}
% Preâmbulo compartilhado: MATF14, Listas de Exercícios 2026
% Usado por ListaNN.tex e GabaritoNN.tex via % Preâmbulo compartilhado: MATF14, Listas de Exercícios 2026
% Usado por ListaNN.tex e GabaritoNN.tex via % Preâmbulo compartilhado: MATF14, Listas de Exercícios 2026
% Usado por ListaNN.tex e GabaritoNN.tex via \input{preamble}

\usepackage[utf8]{inputenc}
\usepackage[T1]{fontenc}
\usepackage[brazil]{babel}
\usepackage{fullpage}
\usepackage{amsmath,amssymb}
\usepackage{enumitem}
\usepackage{xcolor}
\usepackage{fancyhdr}
\usepackage{titlesec}
\usepackage{listings}
\usepackage{hyperref}
\usepackage{booktabs}
\usepackage{graphicx}
\usepackage{tikz}

\definecolor{matf14azul}{HTML}{0D3B54}
\definecolor{matf14dourado}{HTML}{C9971E}

\titleformat{\section}{\normalfont\Large\bfseries\color{matf14azul}}{\thesection}{1em}{}
\titleformat{\subsection}{\normalfont\large\bfseries\color{matf14azul}}{\thesubsection}{1em}{}

\providecommand{\ListaTitulo}{Lista de Exercícios}

\setlength{\headheight}{14pt}
\pagestyle{fancy}
\fancyhf{}
\lhead{\small MATF14: Estatística Econômica I}
\rhead{\small \ListaTitulo}
\cfoot{\thepage}
\renewcommand{\headrulewidth}{0.4pt}

\lstset{
  basicstyle=\ttfamily\small,
  breaklines=true,
  frame=single,
  columns=fullflexible,
  backgroundcolor=\color{gray!8},
  commentstyle=\color{matf14dourado},
  keywordstyle=\color{matf14azul}\bfseries,
  language=R
}

\newcounter{questaocounter}
\newenvironment{questao}[1]{%
  \stepcounter{questaocounter}%
  \bigskip\noindent\textbf{Questão #1.}\ %
}{\par}

\newcommand{\cabecalho}[2]{%
  \begin{center}
    {\Large\bfseries\color{matf14azul} #1}\\[0.3em]
    {\large #2}
  \end{center}
  \vspace{0.5em}
  \noindent\rule{\linewidth}{0.6pt}
  \vspace{0.5em}
}

\newenvironment{solucao}{%
  \par\smallskip\noindent\textit{\color{matf14dourado!80!black}Solução.}\ %
}{\hfill$\blacksquare$\par}


\usepackage[utf8]{inputenc}
\usepackage[T1]{fontenc}
\usepackage[brazil]{babel}
\usepackage{fullpage}
\usepackage{amsmath,amssymb}
\usepackage{enumitem}
\usepackage{xcolor}
\usepackage{fancyhdr}
\usepackage{titlesec}
\usepackage{listings}
\usepackage{hyperref}
\usepackage{booktabs}
\usepackage{graphicx}
\usepackage{tikz}

\definecolor{matf14azul}{HTML}{0D3B54}
\definecolor{matf14dourado}{HTML}{C9971E}

\titleformat{\section}{\normalfont\Large\bfseries\color{matf14azul}}{\thesection}{1em}{}
\titleformat{\subsection}{\normalfont\large\bfseries\color{matf14azul}}{\thesubsection}{1em}{}

\providecommand{\ListaTitulo}{Lista de Exercícios}

\setlength{\headheight}{14pt}
\pagestyle{fancy}
\fancyhf{}
\lhead{\small MATF14: Estatística Econômica I}
\rhead{\small \ListaTitulo}
\cfoot{\thepage}
\renewcommand{\headrulewidth}{0.4pt}

\lstset{
  basicstyle=\ttfamily\small,
  breaklines=true,
  frame=single,
  columns=fullflexible,
  backgroundcolor=\color{gray!8},
  commentstyle=\color{matf14dourado},
  keywordstyle=\color{matf14azul}\bfseries,
  language=R
}

\newcounter{questaocounter}
\newenvironment{questao}[1]{%
  \stepcounter{questaocounter}%
  \bigskip\noindent\textbf{Questão #1.}\ %
}{\par}

\newcommand{\cabecalho}[2]{%
  \begin{center}
    {\Large\bfseries\color{matf14azul} #1}\\[0.3em]
    {\large #2}
  \end{center}
  \vspace{0.5em}
  \noindent\rule{\linewidth}{0.6pt}
  \vspace{0.5em}
}

\newenvironment{solucao}{%
  \par\smallskip\noindent\textit{\color{matf14dourado!80!black}Solução.}\ %
}{\hfill$\blacksquare$\par}


\usepackage[utf8]{inputenc}
\usepackage[T1]{fontenc}
\usepackage[brazil]{babel}
\usepackage{fullpage}
\usepackage{amsmath,amssymb}
\usepackage{enumitem}
\usepackage{xcolor}
\usepackage{fancyhdr}
\usepackage{titlesec}
\usepackage{listings}
\usepackage{hyperref}
\usepackage{booktabs}
\usepackage{graphicx}
\usepackage{tikz}

\definecolor{matf14azul}{HTML}{0D3B54}
\definecolor{matf14dourado}{HTML}{C9971E}

\titleformat{\section}{\normalfont\Large\bfseries\color{matf14azul}}{\thesection}{1em}{}
\titleformat{\subsection}{\normalfont\large\bfseries\color{matf14azul}}{\thesubsection}{1em}{}

\providecommand{\ListaTitulo}{Lista de Exercícios}

\setlength{\headheight}{14pt}
\pagestyle{fancy}
\fancyhf{}
\lhead{\small MATF14: Estatística Econômica I}
\rhead{\small \ListaTitulo}
\cfoot{\thepage}
\renewcommand{\headrulewidth}{0.4pt}

\lstset{
  basicstyle=\ttfamily\small,
  breaklines=true,
  frame=single,
  columns=fullflexible,
  backgroundcolor=\color{gray!8},
  commentstyle=\color{matf14dourado},
  keywordstyle=\color{matf14azul}\bfseries,
  language=R
}

\newcounter{questaocounter}
\newenvironment{questao}[1]{%
  \stepcounter{questaocounter}%
  \bigskip\noindent\textbf{Questão #1.}\ %
}{\par}

\newcommand{\cabecalho}[2]{%
  \begin{center}
    {\Large\bfseries\color{matf14azul} #1}\\[0.3em]
    {\large #2}
  \end{center}
  \vspace{0.5em}
  \noindent\rule{\linewidth}{0.6pt}
  \vspace{0.5em}
}

\newenvironment{solucao}{%
  \par\smallskip\noindent\textit{\color{matf14dourado!80!black}Solução.}\ %
}{\hfill$\blacksquare$\par}


\begin{document}

\cabecalho{Lista de Exercícios 06}{Probabilidade condicional e independência estatística (Cap. 2, Aulas 11--12)}

\noindent \textit{Justifique todas as respostas. As resoluções são discutidas em sala e nos horários de atendimento; não são publicadas neste material.}

\begin{questao}{1} \textbf{(Condicional: inadimplência por faixa de renda)}

Uma financeira tem 1.000 clientes classificados por faixa de renda e situação de pagamento:

\begin{center}
\begin{tabular}{l|cc|c}
\toprule
 & Em dia & Inadimplente & Total \\
\midrule
Renda baixa  & 320 & 80  & 400 \\
Renda média  & 380 & 20  & 400 \\
Renda alta   & 195 & 5   & 200 \\
\midrule
Total        & 895 & 105 & 1000\\
\bottomrule
\end{tabular}
\end{center}

\begin{enumerate}[label=(\alph*)]
  \item Calcule $P(\text{inadimplente}\mid\text{renda baixa})$.
  \item Calcule $P(\text{inadimplente}\mid\text{renda alta})$.
  \item A situação de pagamento é independente da faixa de renda nesta base? Justifique
    comparando as probabilidades condicionais entre si e com $P(\text{inadimplente})$ geral.
\end{enumerate}
\end{questao}

\begin{questao}{2} \textbf{(Regra do produto: sorteio sem reposição)}

Um comitê de auditoria tem 15 processos para revisar, dos quais 4 têm irregularidades. Dois
processos são sorteados, sem reposição, para revisão detalhada nesta semana.

\begin{enumerate}[label=(\alph*)]
  \item Qual a probabilidade de que os dois processos sorteados tenham irregularidades?
  \item Qual a probabilidade de que \emph{nenhum} dos dois tenha irregularidade?
  \item Refaça o item (a) supondo, em vez disso, que o sorteio fosse \emph{com} reposição (o
    primeiro processo sorteado volta para o conjunto antes do segundo sorteio). Compare com o
    resultado original e explique a diferença.
\end{enumerate}
\end{questao}

\begin{questao}{3} \textbf{(Independência: verificação numérica)}

Em uma pesquisa com 200 pequenas empresas, 120 são do setor de comércio e 60 relataram ter
acessado alguma linha de crédito subsidiado no último ano. Dessas 60, 36 são do setor de comércio.

\begin{enumerate}[label=(\alph*)]
  \item Calcule $P(\text{comércio})$, $P(\text{acessou crédito})$ e
    $P(\text{comércio}\cap\text{acessou crédito})$.
  \item Verifique se $P(\text{comércio}\cap\text{acessou crédito}) = P(\text{comércio}) \times
    P(\text{acessou crédito})$. Setor e acesso a crédito são independentes nesta amostra?
\end{enumerate}
\end{questao}

\begin{questao}{4} \textbf{(Mutuamente exclusivos vs. independentes: discussão numérica)}

Considere dois eventos $A$ e $B$, com $P(A)=0{,}4$ e $P(B)=0{,}5$.

\begin{enumerate}[label=(\alph*)]
  \item Se $A$ e $B$ forem mutuamente exclusivos, qual é $P(A\cap B)$? E $P(A\mid B)$?
  \item Se $A$ e $B$ forem independentes, qual é $P(A\cap B)$? E $P(A\mid B)$?
  \item É possível que $A$ e $B$, com essas probabilidades individuais, sejam \emph{ao mesmo
    tempo} mutuamente exclusivos e independentes? (Dica: compare os dois valores de $P(A\cap B)$
    encontrados acima e o que a definição de independência exigiria de $P(A\mid B)$ se
    $P(A\cap B)=0$.)
\end{enumerate}
\end{questao}

\begin{questao}{5} \textbf{(Dedução: independência é simétrica)}

Suponha $P(A\cap B) = P(A)P(B)$ (definição de independência), com $P(A)>0$ e $P(B)>0$.

\begin{enumerate}[label=(\alph*)]
  \item Mostre que $P(A\mid B) = P(A)$.
  \item Mostre que $P(B\mid A) = P(B)$ (ou seja, a independência é uma relação simétrica: se $A$
    é independente de $B$, então $B$ também é independente de $A$, não é preciso verificar as
    duas direções separadamente na prática).
\end{enumerate}
\end{questao}

\bigskip
\noindent\textit{A questão a seguir não tem resposta única e não é cobrada em prova no mesmo
formato, é para discussão em grupo ou por escrito, antes da próxima aula.}

\begin{questao}{6 (discussão)} \textbf{(Sem resposta única)}

Um analista de crédito assume, por conveniência de cálculo, que a inadimplência de diferentes
tomadores de um mesmo banco é independente entre si. Escreva um parágrafo explicando por que essa
suposição pode ser razoável em tempos normais, mas perigosa durante uma crise econômica
generalizada, e que tipo de risco para o banco fica escondido quando essa suposição é usada sem
qualificação.
\end{questao}


\bigskip
\section*{Exercícios complementares}

\begin{enumerate}[label=\textbf{C\arabic*.}, leftmargin=*]
\item Um restaurante popular serve salada ou um prato à base de carne. 75\% dos fregueses são homens, 10\% dos homens preferem salada e 40\% das mulheres preferem carne. Um freguês é sorteado; sejam $H$ (homem), $M$ (mulher), $S$ (prefere salada) e $C$ (prefere carne). Calcule:
\begin{enumerate}[label=(\alph*)]
  \item $P(H)$, $P(S \mid H)$ e $P(C \mid H)$;
  \item $P(S \cup H)$ e $P(S \cap H)$;
  \item $P(M \mid S)$.
\end{enumerate}
\item Os 250 estudantes de uma faculdade, por sexo e área:
\begin{center}
\begin{tabular}{l|ccc}
\toprule
 & Biológicas & Exatas & Humanas \\
\midrule
Masculino & 42 & 30 & 58 \\
Feminino & 38 & 12 & 70 \\
\bottomrule
\end{tabular}
\end{center}
Um estudante é sorteado. Qual a probabilidade de ele ser
\begin{enumerate}[label=(\alph*)]
  \item do sexo feminino e da área de Humanas?
  \item do sexo masculino e não ser de Biológicas?
  \item do sexo feminino, dado que é de Humanas?
  \item Os eventos \emph{feminino} e \emph{Humanas} são independentes?
\end{enumerate}
\item Uma comissão de três moradores é sorteada numa comunidade com 50\% de moradores brancos, 30\% negros e 20\% latinos; o primeiro sorteado é o líder. Suponha sorteios independentes. Qual a probabilidade de
\begin{enumerate}[label=(\alph*)]
  \item o líder ser negro?
  \item o líder ser branco ou latino?
  \item os três membros serem negros?
\end{enumerate}
\item Suponha que 16\% dos presos em prisões estaduais sofram de alguma doença psiquiátrica. Qual a probabilidade de
\begin{enumerate}[label=(\alph*)]
  \item um preso escolhido ao acaso não sofrer de doença psiquiátrica?
  \item nenhum de três presos independentes sofrer de doença psiquiátrica?
\end{enumerate}
\item A probabilidade de a pessoa $A$ resolver um problema é $2/3$ e a de $B$ é $3/4$. Se tentarem independentemente, qual a probabilidade de o problema ser resolvido?
\item Uma carta é sorteada de um baralho comum de 52 cartas. Seja $R$ o evento \emph{rei} e $E$ o evento \emph{espadas}. $R$ e $E$ são independentes? São disjuntos?

\end{enumerate}

\bigskip
\needspace{12\baselineskip}
\section*{Aprofundamento: novos tópicos e dados reais}
\noindent\textit{Exercícios sobre os tópicos acrescentados ao Capítulo 2 do livro e às aulas.}

\begin{enumerate}[label=\textbf{A\arabic*.}, leftmargin=*]

\item \textbf{(Probabilidade geométrica.)}
\begin{enumerate}[label=(\alph*)]
  \item Dois sócios chegam a uma reunião em instantes ao acaso entre 9h e 10h, independentemente, e cada um espera 10
    minutos pelo outro. Qual a probabilidade de se encontrarem?
  \item Um ônibus passa a cada 12 minutos. Um passageiro chega num instante ao acaso. Qual a probabilidade de esperar
    mais de 5 minutos? E de esperar exatamente 5 minutos?
  \item O preço máximo que um comprador aceita pagar, $b$, e o mínimo que o vendedor aceita receber, $s$, são pontos ao
    acaso em $[0, 1]$, independentes. Qual a probabilidade de o ganho de troca $b - s$ ser pelo menos $0{,}3$?
\end{enumerate}

\item \textbf{(Tabela de contingência com dados reais.)} Pessoas de 14 anos ou mais no Brasil, 2º trimestre de 2026 (mil pessoas;
IBGE, PNAD Contínua):
\begin{center}
\begin{tabular}{l|ccc|c}
\toprule
 & Ocupadas & Desocupadas & Fora da força & Total \\
\midrule
Homens & 58.293 & 2.803 & 23.465 & 84.561 \\
Mulheres & 44.765 & 3.059 & 43.016 & 90.840 \\
\midrule
Total & 103.058 & 5.862 & 66.481 & 175.401 \\
\bottomrule
\end{tabular}
\end{center}
{\footnotesize Totais obtidos somando as linhas; diferem em até 1 mil das estimativas publicadas por causa do arredondamento.}

\begin{enumerate}[label=(\alph*)]
  \item Calcule $P(\text{mulher} \cap \text{fora da força})$, $P(\text{fora da força})$, $P(\text{fora da força} \mid \text{mulher})$ e
    $P(\text{mulher} \mid \text{fora da força})$. Quais são conjuntas, marginais e condicionais?
  \item A \emph{taxa} de desocupação é $P(\text{desocupado} \mid \text{na força de trabalho})$ e o \emph{nível} de desocupação é
    $P(\text{desocupado} \mid \text{14 anos ou mais})$. Calcule os dois para o Brasil.
  \item Sob independência entre sexo e condição, quantas mulheres (em mil) estariam fora da força de trabalho? Compare com o observado.
\end{enumerate}

\item \textbf{(Paradoxo de Simpson.)} Duas agências de um banco receberam pedidos de crédito de pequenas e médias empresas:
\begin{center}
\begin{tabular}{l|cc|cc}
\toprule
 & \multicolumn{2}{c|}{Agência X} & \multicolumn{2}{c}{Agência Y} \\
 & pedidos & aprovados & pedidos & aprovados \\
\midrule
Pequenas & 100 & 40 & 400 & 180 \\
Médias & 400 & 300 & 100 & 80 \\
\bottomrule
\end{tabular}
\end{center}
Calcule as taxas de aprovação por porte e no total em cada agência. Qual agência é ``mais generosa''? Explique a aparente contradição.

\item \textbf{(Risco relativo e razão de chances.)} Entre os ocupados no 2º trimestre de 2026, 22.717 mil de 58.293 mil homens e
15.862 mil de 44.765 mil mulheres estavam em situação de informalidade (IBGE, PNAD Contínua).
\begin{enumerate}[label=(\alph*)]
  \item Calcule a taxa de informalidade de cada sexo, a diferença de riscos, o risco relativo e a razão de chances (homens em relação às mulheres).
  \item Por que a razão de chances é maior que o risco relativo? Em que situação as duas medidas ficariam próximas?
  \item Um pesquisador sorteou 1.000 trabalhadores informais e 1.000 formais e comparou a proporção de homens. Qual das duas medidas
    ele consegue estimar corretamente com esse desenho? Por quê?
\end{enumerate}

\item \textbf{(Independência e complementares.)} Sejam $A$ e $B$ independentes com $P(A) = 0{,}3$ e $P(B) = 0{,}5$. Calcule $P(A \cup B)$ e
verifique numericamente que $A^c$ e $B$ também são independentes.

\end{enumerate}

\end{document}
